\section{预告普查结论：16只不是同一种高增长}

54只核心/卫星池中，16只在截至7月11日捕获的H1预告表中披露业绩，其中15只预增、1只扭亏。判断顺序是：先看H1利润和隐含Q2，再看扣非与一次性，最后比较H1预告EPS与最新券商全年EPS。若H1已经超过旧全年分母，旧PE立即失效，必须进入预告后重估池。

\small
\begin{longtable}{L{1.35cm}L{1.75cm}R{1.05cm}R{1.05cm}R{1.10cm}R{1.15cm}R{1.10cm}L{2.45cm}}
\toprule
\textbf{代码} & \textbf{公司} & \textbf{H1净利} & \textbf{隐含Q2} & \textbf{Q2/Q1} & \textbf{H1 EPS} & \textbf{占旧全年} & \textbf{质量结论} \\
\midrule
\endfirsthead
\toprule
\textbf{代码} & \textbf{公司} & \textbf{H1净利} & \textbf{隐含Q2} & \textbf{Q2/Q1} & \textbf{H1 EPS} & \textbf{占旧全年} & \textbf{质量结论} \\
\midrule
\endhead
601138 & 工业富联 & 239.00 & 133.05 & 25.6\% & 1.20 & 42.4\% & 主营真兑现 \\
000977 & 浪潮信息 & 28.50 & 22.45 & 271.1\% & 1.94 & 76.7\% & 主营真兑现 \\
301165 & 锐捷网络 & 6.75 & 5.52 & 349.1\% & 0.85 & 45.4\% & 主营真兑现 \\
600601 & 方正科技 & 5.70 & 3.38 & 45.9\% & 0.12 & 111.8\% & 主营真兑现 \\
001339 & 智微智能 & 3.84 & 2.74 & 151.1\% & 1.51 & 60.9\% & 主营真兑现 \\
301308 & 江波龙 & 101.00 & 62.38 & 61.5\% & 24.08 & 91.0\% & 周期/结构真兑现 \\
300475 & 香农芯创 & 37.50 & 24.23 & 82.5\% & 8.08 & 342.4\% & 周期/结构真兑现 \\
002407 & 多氟多 & 5.05 & 1.29 & -65.7\% & 0.43 & 119.5\% & 周期/结构真兑现 \\
000636 & 风华高科 & 2.85 & 1.96 & 121.3\% & 0.26 & 61.3\% & 周期/结构真兑现 \\
600685 & 中船防务 & 8.40 & 4.44 & 12.1\% & 0.59 & 55.5\% & 主营改善+投资收益 \\
002842 & 翔鹭钨业 & 5.50 & 2.99 & 19.4\% & 1.69 & --- & 周期/价格兑现 \\
000938 & 紫光股份 & 21.15 & 13.27 & 68.4\% & 0.74 & 87.0\% & 主营+结构变化+一次性 \\
603986 & 兆易创新 & 69.00 & 54.39 & 272.2\% & 10.34 & 123.8\% & 主营+公允价值收益 \\
002396 & 星网锐捷 & 3.70 & 3.33 & 794.5\% & 0.64 & 52.6\% & 主营+一次性 \\
001314 & 亿道信息 & 1.96 & 1.53 & 256.6\% & 1.37 & 82.5\% & 主营+一次性 \\
002829 & 星网宇达 & 1.05 & 1.32 & -597.5\% & 0.51 & 39.6\% & 非经常损益主导 \\
\bottomrule
\end{longtable}
\normalsize
\sourcenote{16份公司正式业绩预告PDF、2026Q1财务包、最新券商预测元数据；利润单位为亿元。}

\section{第一组：主营真兑现，可进入正式模型或回撤观察}

\textbf{工业富联}的AI服务器收入、800G以上交换机出货和Q2利润均有直接证据，且已有华泰/华创预告后预测，因此进入正式估值。\textbf{浪潮信息}H1净利中值28.5亿元、隐含Q2约22.45亿元，Q2较Q1增长271\%；群益预告同日给出56.67亿元2026E净利和90元目标，中邮预告后EPS 2.53元。但7月10日现价89.52元几乎到达90元目标，结论是业绩兑现、等待回撤，而非继续追高。

\textbf{锐捷网络}隐含Q2净利5.52亿元、较Q1增长349\%，非经常损益占H1约2.2\%，数据中心交换机是清晰主营驱动；但20日涨幅超过80\%、一年位置约92\%，预告后开源EPS 1.87元对应现价估值仍高，只进入回撤观察。\textbf{方正科技}H1净利5.7亿元、扣非方向与归母一致，高端PCB产品结构改善；但H1 EPS约0.12元已超过旧全年0.11元，必须重建全年分母。\textbf{智微智能}H1净利3.84亿元、Q2明显加速，但股价已在一年高位，先等待预告后研报与现金流。

\section{第二组：周期与产品结构兑现，但不能把高景气线性外推}

\textbf{江波龙}H1净利中值101亿元，H1 EPS约24.08元，已完成旧国信全年EPS26.47元的91\%；本报告使用预告后自行重建的200亿元全年净利，而非旧预测。\textbf{香农芯创}H1 EPS约8.08元，是旧全年2.36元的342\%，旧PE完全失效；企业级存储和海普存储商业化是真增长，但分销毛利、存货和现金流必须重做。

\textbf{风华高科}Q2隐含利润较Q1增长121\%，MLCC、电阻、电感量价和降本共同驱动；现价和周期位置要求持续改善。\textbf{多氟多}H1 EPS约0.43元已超过旧全年0.36元，但隐含Q2利润较Q1下降65.7\%，不能只看同比高增；需验证六氟磷酸锂价格、销量和现金流。\textbf{翔鹭钨业}受益钨价传导，Q2继续增长，但缺少有效2026E券商分母，高钨价能否顺利向下游转嫁是核心风险。

\section{第三组：主营改善中混入结构性或一次性收益}

\textbf{紫光股份}H1净利中值21.15亿元，AI/ICT增长和新华三持股提升均为正面，但约2.5--3.5亿元非经常收益不可持续；H1 EPS约0.74元已达到旧全年0.85元的87\%，旧模型必须同时更新持股比例、利息费用和一次性收益。\textbf{兆易创新}H1净利69亿元、隐含Q2约54.39亿元，但包含证券公允价值收益；H1 EPS约10.34元已超过旧全年8.35元，重估必须使用扣非利润。

\textbf{星网锐捷}数据中心交换机是真驱动，但约1.0--1.2亿元非经常收益占H1中值约29.7\%；\textbf{亿道信息}约0.86亿元投资收益占H1中值约43.9\%，均需把主营与一次性分开。\textbf{中船防务}造船订单、生产效率和毛利改善是真实主线，但联营企业与分红收益也推高利润，因此进入预告后重估池而非直接给目标。

\section{明确排除：星网宇达}

星网宇达H1预计扭亏至约1.05亿元，但扣非仍亏，约1.4亿元投资收益是扭亏主因，一次性占报告利润中值超过100\%。在主营收入下降、扣非未转正前，该预告不能获得盈利信用，明确排除出核心和重估池。

\begin{riskbox}[预告使用边界]
业绩预告未经审计；H1利润不是全年利润的简单一半。未披露扣非、现金流、存货和客户订单的标的，只能得到方向性信用。预告后没有新研报或公司级模型时，本报告只给“重估池/等待模型刷新”，不伪造点目标。
\end{riskbox}
